% Options for packages loaded elsewhere
% Options for packages loaded elsewhere
\PassOptionsToPackage{unicode}{hyperref}
\PassOptionsToPackage{hyphens}{url}
\PassOptionsToPackage{dvipsnames,svgnames,x11names}{xcolor}
%
\documentclass[
  letterpaper,
  DIV=11,
  numbers=noendperiod]{scrartcl}
\usepackage{xcolor}
\usepackage{amsmath,amssymb}
\setcounter{secnumdepth}{-\maxdimen} % remove section numbering
\usepackage{iftex}
\ifPDFTeX
  \usepackage[T1]{fontenc}
  \usepackage[utf8]{inputenc}
  \usepackage{textcomp} % provide euro and other symbols
\else % if luatex or xetex
  \usepackage{unicode-math} % this also loads fontspec
  \defaultfontfeatures{Scale=MatchLowercase}
  \defaultfontfeatures[\rmfamily]{Ligatures=TeX,Scale=1}
\fi
\usepackage{lmodern}
\ifPDFTeX\else
  % xetex/luatex font selection
\fi
% Use upquote if available, for straight quotes in verbatim environments
\IfFileExists{upquote.sty}{\usepackage{upquote}}{}
\IfFileExists{microtype.sty}{% use microtype if available
  \usepackage[]{microtype}
  \UseMicrotypeSet[protrusion]{basicmath} % disable protrusion for tt fonts
}{}
\makeatletter
\@ifundefined{KOMAClassName}{% if non-KOMA class
  \IfFileExists{parskip.sty}{%
    \usepackage{parskip}
  }{% else
    \setlength{\parindent}{0pt}
    \setlength{\parskip}{6pt plus 2pt minus 1pt}}
}{% if KOMA class
  \KOMAoptions{parskip=half}}
\makeatother
% Make \paragraph and \subparagraph free-standing
\makeatletter
\ifx\paragraph\undefined\else
  \let\oldparagraph\paragraph
  \renewcommand{\paragraph}{
    \@ifstar
      \xxxParagraphStar
      \xxxParagraphNoStar
  }
  \newcommand{\xxxParagraphStar}[1]{\oldparagraph*{#1}\mbox{}}
  \newcommand{\xxxParagraphNoStar}[1]{\oldparagraph{#1}\mbox{}}
\fi
\ifx\subparagraph\undefined\else
  \let\oldsubparagraph\subparagraph
  \renewcommand{\subparagraph}{
    \@ifstar
      \xxxSubParagraphStar
      \xxxSubParagraphNoStar
  }
  \newcommand{\xxxSubParagraphStar}[1]{\oldsubparagraph*{#1}\mbox{}}
  \newcommand{\xxxSubParagraphNoStar}[1]{\oldsubparagraph{#1}\mbox{}}
\fi
\makeatother


\usepackage{longtable,booktabs,array}
\usepackage{calc} % for calculating minipage widths
% Correct order of tables after \paragraph or \subparagraph
\usepackage{etoolbox}
\makeatletter
\patchcmd\longtable{\par}{\if@noskipsec\mbox{}\fi\par}{}{}
\makeatother
% Allow footnotes in longtable head/foot
\IfFileExists{footnotehyper.sty}{\usepackage{footnotehyper}}{\usepackage{footnote}}
\makesavenoteenv{longtable}
\usepackage{graphicx}
\makeatletter
\newsavebox\pandoc@box
\newcommand*\pandocbounded[1]{% scales image to fit in text height/width
  \sbox\pandoc@box{#1}%
  \Gscale@div\@tempa{\textheight}{\dimexpr\ht\pandoc@box+\dp\pandoc@box\relax}%
  \Gscale@div\@tempb{\linewidth}{\wd\pandoc@box}%
  \ifdim\@tempb\p@<\@tempa\p@\let\@tempa\@tempb\fi% select the smaller of both
  \ifdim\@tempa\p@<\p@\scalebox{\@tempa}{\usebox\pandoc@box}%
  \else\usebox{\pandoc@box}%
  \fi%
}
% Set default figure placement to htbp
\def\fps@figure{htbp}
\makeatother


% definitions for citeproc citations
\NewDocumentCommand\citeproctext{}{}
\NewDocumentCommand\citeproc{mm}{%
  \begingroup\def\citeproctext{#2}\cite{#1}\endgroup}
\makeatletter
 % allow citations to break across lines
 \let\@cite@ofmt\@firstofone
 % avoid brackets around text for \cite:
 \def\@biblabel#1{}
 \def\@cite#1#2{{#1\if@tempswa , #2\fi}}
\makeatother
\newlength{\cslhangindent}
\setlength{\cslhangindent}{1.5em}
\newlength{\csllabelwidth}
\setlength{\csllabelwidth}{3em}
\newenvironment{CSLReferences}[2] % #1 hanging-indent, #2 entry-spacing
 {\begin{list}{}{%
  \setlength{\itemindent}{0pt}
  \setlength{\leftmargin}{0pt}
  \setlength{\parsep}{0pt}
  % turn on hanging indent if param 1 is 1
  \ifodd #1
   \setlength{\leftmargin}{\cslhangindent}
   \setlength{\itemindent}{-1\cslhangindent}
  \fi
  % set entry spacing
  \setlength{\itemsep}{#2\baselineskip}}}
 {\end{list}}
\usepackage{calc}
\newcommand{\CSLBlock}[1]{\hfill\break\parbox[t]{\linewidth}{\strut\ignorespaces#1\strut}}
\newcommand{\CSLLeftMargin}[1]{\parbox[t]{\csllabelwidth}{\strut#1\strut}}
\newcommand{\CSLRightInline}[1]{\parbox[t]{\linewidth - \csllabelwidth}{\strut#1\strut}}
\newcommand{\CSLIndent}[1]{\hspace{\cslhangindent}#1}



\setlength{\emergencystretch}{3em} % prevent overfull lines

\providecommand{\tightlist}{%
  \setlength{\itemsep}{0pt}\setlength{\parskip}{0pt}}



 


\usepackage{xr}
\externaldocument{ab_supp}
\newcommand{\sfig}[1]{Supplementary Figure~\ref{#1}}
\newcommand{\ssec}[1]{Supplementary Section~\ref{#1}}
\usepackage{lineno}
\KOMAoption{captions}{tableheading}
\makeatletter
\@ifpackageloaded{caption}{}{\usepackage{caption}}
\AtBeginDocument{%
\ifdefined\contentsname
  \renewcommand*\contentsname{Table of contents}
\else
  \newcommand\contentsname{Table of contents}
\fi
\ifdefined\listfigurename
  \renewcommand*\listfigurename{List of Figures}
\else
  \newcommand\listfigurename{List of Figures}
\fi
\ifdefined\listtablename
  \renewcommand*\listtablename{List of Tables}
\else
  \newcommand\listtablename{List of Tables}
\fi
\ifdefined\figurename
  \renewcommand*\figurename{Figure}
\else
  \newcommand\figurename{Figure}
\fi
\ifdefined\tablename
  \renewcommand*\tablename{Table}
\else
  \newcommand\tablename{Table}
\fi
}
\@ifpackageloaded{float}{}{\usepackage{float}}
\floatstyle{ruled}
\@ifundefined{c@chapter}{\newfloat{codelisting}{h}{lop}}{\newfloat{codelisting}{h}{lop}[chapter]}
\floatname{codelisting}{Listing}
\newcommand*\listoflistings{\listof{codelisting}{List of Listings}}
\makeatother
\makeatletter
\makeatother
\makeatletter
\@ifpackageloaded{caption}{}{\usepackage{caption}}
\@ifpackageloaded{subcaption}{}{\usepackage{subcaption}}
\makeatother
\usepackage{bookmark}
\IfFileExists{xurl.sty}{\usepackage{xurl}}{} % add URL line breaks if available
\urlstyle{same}
\hypersetup{
  pdftitle={Spatial aggregation confounds null model-based inference of species association networks},
  pdfauthor={Andrew J. Rominger},
  colorlinks=true,
  linkcolor={blue},
  filecolor={Maroon},
  citecolor={Blue},
  urlcolor={Blue},
  pdfcreator={LaTeX via pandoc}}


\title{Spatial aggregation confounds null model-based inference of
species association networks}
\author{Andrew J. Rominger}
\date{}
\begin{document}
\maketitle


\linenumbers

\subsection{Abstract}\label{abstract}

\subsection{Introduction}\label{introduction}

Species interactions are fundamental to the evolutionary origins of
biodiversity and the maintenance of that diversity in ecological
assemblages (Hutchinson 1961, Holt and Lawton 1994, Chesson 2000, Proulx
et al. 2005, Bascompte and Jordano 2007, Levine and HilleRisLambers
2009, Allesina and Tang 2012, Johnson and Bronstein 2019). Network
theory has come to be a central paradigm in ecology and evolution (May
1972, Williams and Martinez 2000, Dunne et al. 2002, Proulx et al. 2005,
Bascompte and Jordano 2007), and so, not surprisingly, these
interactions are often cast as ecological networks where species are
nodes (or vertices) and network links (or edges) connect species that
interact.

Abundance or presence/absence data are among the most prevalent and
accessible ecological data, and have consequentially been frequently
leveraged (Morueta-Holme et al. 2016, Calatayud et al. 2019) in an
attempt to infer species interactions (or more accurately associations),
beginning at least in the work of Diamond (1975). The specific data used
for constructing association networks are often site by species matrices
where each cell in the matrix records the abundance (or
presence/absence) of a given species at a given site (i.e.~sampling
location). For simplicity I will refer to both abundance and
presence/absence data as simply ``abundance data'' and focus on true
abundance data with the assumed implication that methodological
limitations in the case of true abundances will be even more confounding
for presence/absence scenarios (Ulrich and Gotelli 2010).

Two salient insights from decades of debate about using abundance data
to infer ecological processes are: 1) interactions in the strict sense
cannot be inferred, but \emph{species associations} might be recoverable
and these \emph{associations} are assumed to map onto a combination of
species-species and species-environment interactions (Morueta-Holme et
al. 2016, Calatayud et al. 2019); and 2) robust null models are
necessary to identify which patterns in the data are likely to arise by
chance and which are likely to carry biological signal (Ulrich and
Gotelli 2010).

Null models seek to generate plausible data, and patterns in data, while
stripping away the actual relationships between species in the data
(Ulrich and Gotelli 2010). Most often a null model is achieved by
repeatedly permuting the original data such that any real associations
are removed because which species co-occur has been randomized by
permutation (Ulrich and Gotelli 2010). The collection of all such
permutations represents a null distribution. These permutation
approaches also seek to preserve key aspects of the original data such
as maintaining the observed abundances of each species and/or
maintaining the observed total number of individuals at each site (Ladau
2017). These constraints are imposed in an attempt to control the rate
of incorrect inference of the presence or absence of network links
(Ladau 2008, 2017). Other methods for inferring association networks do
exist, especially model-based inference applied in microbial ecology
(e.g. Zhou et al. 2010, Kurtz et al. 2015), but are less frequently
employed in community ecology (but see, e.g., Harris 2016, Losapio et
al. 2021).

Here I raise a previously unreported technical issue with the approach
of inferring association networks from abundance data using
permutational null models. The issue arises from the near universal
spatial aggregation of individuals within species (Legendre 1993, Condit
et al. 2000, Lennon 2000, Dale and Fortin 2002, McGill and Collins 2003,
Engen et al. 2008, Zillio and He 2010, Harte 2011, Connolly et al.
2017). Using simulation experiments and analysis of real data (from
Calatayud et al. 2019), I show that intra-specific spatial aggregation
(ISSA) causes exceedingly high type I error rates in permutational null
model approaches to inferring association networks, i.e., incorrectly
detecting associations when there are none. ISSA in itself is not
evidence of species association---it can arise from a large number of
processes including simple demographic drift (Kendall 1949, Hubbell
2001) in which species do not meaningfully interact with each other or
shared environments.

In the absence of ISSA, abundances of a species across sites should
follow a Poisson distribution (Diggle 2003). Thus ISSA is often
quantified as the deviation of observed data from a Poisson distribution
(Diggle 2003). Conversely, the negative binomial distribution can model
spatially aggregated intra-specific abundances across sites (Bliss and
Fisher 1953, Taylor 1961, Engen et al. 2008, Connolly et al. 2017).
Critically, the negative binomial arises very naturally from ecological
processes as the distribution of population sizes in a
birth-death-immigration model (Kendall 1949). If the abundance of a
species at a site is determined by local birth and death with
supplemental individuals arriving by immigration from a source pool,
then the probability distribution of repeatedly sampled abundances will
follow a negative binomial distribution (Kendall 1949). In this way,
simple demographic drift (birth, death, and immigration that are purely
stochastic) is expected to produce intra-specific abundances that are
negative binomial-distributed and spatially aggregated. Thus ISSA is
not, by itself, evidence of species association.

The technical issue raised here is specific to the widely used family of
methods that infer association networks by comparing an observed test
statistic to a null distribution generated by Monte Carlo permutation of
the site-by-species matrix. I do not test whether the same issue affects
methods that do not rely on this style of null model, most notably joint
species distribution models, which explicitly condition on measured or
latent environmental covariates and site-level random effects (e.g. Zhou
et al. 2010), or the compositional correlation methods common in
microbial ecology (e.g. Kurtz et al. 2015). It is not obvious that the
mechanism described here transfers directly to these other methods, and
I do not want to imply, by omission, that it does. However, any effects
of ISSA on joint species distribution modeling and compositional
correlation methods should be assessed. Here I focus solely
permutational null models because of their foundational status in the
field (\textbf{cite?}) and their continued use in inferring association
networks.

\subsection{Methods}\label{methods}

\subsubsection{Reproducibility and data}\label{reproducibility-and-data}

A custom R package, \emph{SpatAggNet} accompanies this paper and can be
used to reproduce all results. ssec\{sec-setup\} details how to access
and install \emph{SpatAggNet}. All analysis scripts and data analyzed
are available via the package.

I use a compiled data set of site-by-species abundance matrices
assembled by Calatayud et al. (2019). The data are included with the
\emph{SpatAggNet} package as \texttt{abund\_mats} and comprise
\texttt{\textbar{}r\ length(abund\_mats)} site-by-species matrices.
These compiled data are ideal for investigating the impact of ISSA on
association network inference because they were used by Calatayud et al.
(2019) for the very purpose of network inference and also because they
represent a spatially and taxonomically extensive collection of
quality-controlled data.

\subsubsection{Quantifying spatial
aggregation}\label{quantifying-spatial-aggregation}

The spatial species abundance distribution (SSAD) is the distribution of
a single species' replicate abundances across sites. I fit both negative
binomial and Poisson distributions to the SSADs of each species in each
site-by-species matrix using maximum likelihood. The \texttt{nb\_fit()}
function from \emph{SpatAggNet} fits both distributions to every species
in each site-by-species matrix and returns, for each species from each
matrix:

\begin{itemize}
\tightlist
\item
  the negative binomial dispersion parameter (\texttt{size}, i.e.~\(k\))
  and mean (\texttt{mu})
\item
  the log-likelihoods of negative binomial and Poisson models
\item
  their AIC difference (\texttt{delta\_AIC}),
\item
  a \(\chi^2\)-distributed goodness-of-fit statistic for the negative
  binomial {[}\texttt{zsq}; Etienne (2007){]}
\item
  the species' total abundance and occupancy.
\end{itemize}

\subsubsection{Analysis of empirical
data}\label{analysis-of-empirical-data}

\begin{itemize}
\tightlist
\item
  purpose: show the issue and that nb is the reason
\item
  data from rareplusminus
\item
  get sad and ssad from those data

  \begin{itemize}
  \tightlist
  \item
    pois
  \item
    nb
  \item
    model selection of pois and nb
  \end{itemize}
\item
  look at network stats for real data versus data simulated from sad and
  ssad

  \begin{itemize}
  \tightlist
  \item
    pois: doesn't look like real data
  \item
    nb: looks close to real data
  \end{itemize}
\item
  report on rareplusminus stuff
\end{itemize}

\subsubsection{Simulation experiment}\label{simulation-experiment}

\begin{itemize}
\tightlist
\item
  purpose: show the issue persists across different methods
\item
  start with data simulated from sad and ssad

  \begin{itemize}
  \tightlist
  \item
    only need to do nb
  \end{itemize}
\item
  make networks based on diff perm methods and diff distance metrics
\end{itemize}

\subsubsection{Behavior of negative binomial random variables subjected
to
permutation}\label{behavior-of-negative-binomial-random-variables-subjected-to-permutation}

\begin{itemize}
\tightlist
\item
  show that k is distorted
\end{itemize}

\subsection{Results}\label{results}

\subsubsection{Analysis of empirical
data}\label{analysis-of-empirical-data-1}

\subsubsection{Simulation experiment}\label{simulation-experiment-1}

\subsection{Discussion}\label{discussion}

\subsection*{References}\label{references}
\addcontentsline{toc}{subsection}{References}

\phantomsection\label{refs}
\begin{CSLReferences}{1}{0}
\bibitem[\citeproctext]{ref-allesina2012}
Allesina, S., and S. Tang. 2012. Stability criteria for complex
ecosystems. Nature 483:205--208.

\bibitem[\citeproctext]{ref-bascompte2007}
Bascompte, J., and P. Jordano. 2007. Plant-animal mutualistic networks:
The architecture of biodiversity. Annu. Rev. Ecol. Evol. Syst.
38:567--593.

\bibitem[\citeproctext]{ref-bliss1953}
Bliss, C. I., and R. A. Fisher. 1953. Fitting the negative binomial
distribution to biological data. Biometrics 9:176--200.

\bibitem[\citeproctext]{ref-calatayud2019}
Calatayud, J., E. Andivia, A. Escudero, C. J. Melián, R.
Bernardo-Madrid, M. Stoffel, C. Aponte, N. G. Medina, R. Molina-Venegas,
X. Arnan, M. Rosvall, M. Neuman, J. A. Noriega, F. Alves-Martins, I.
Draper, A. Luzuriaga, J. A. Ballesteros-Cánovas, C. Morales-Molino, P.
Ferrandis, A. Herrero, L. Pataro, L. Juen, A. Cea, and J.
Madrigal-González. 2019. Positive associations among rare species and
their persistence in ecological assemblages. Nat Ecol Evol.

\bibitem[\citeproctext]{ref-chesson2000}
Chesson, P. 2000. Mechanisms of maintenance of species diversity. Annual
review of Ecology and Systematics 31:343--366.

\bibitem[\citeproctext]{ref-condit2000}
Condit, R., P. S. Ashton, P. Baker, S. Bunyavejchewin, S. Gunatilleke,
N. Gunatilleke, S. P. Hubbell, R. B. Foster, A. Itoh, J. V. LaFrankie,
and others. 2000. Spatial patterns in the distribution of tropical tree
species. Science 288:1414--1418.

\bibitem[\citeproctext]{ref-connolly2017}
Connolly, S. R., T. P. Hughes, and D. R. Bellwood. 2017. A unified model
explains commonness and rarity on coral reefs. Ecology letters
20:477--486.

\bibitem[\citeproctext]{ref-dale2002}
Dale, M. R., and M.-J. Fortin. 2002. Spatial autocorrelation and
statistical tests in ecology. Ecoscience 9:162--167.

\bibitem[\citeproctext]{ref-diamond1975}
Diamond, J. M. 1975. Assembly of species communities. Ecology and
evolution of communities:342--444.

\bibitem[\citeproctext]{ref-diggle}
Diggle, Peter. 2003. Statistical analysis of spatial point patterns. 2nd
ed. Arnold, London.

\bibitem[\citeproctext]{ref-dunne2002}
Dunne, J. A., R. J. Williams, and N. D. Martinez. 2002. Network
structure and biodiversity loss in food webs: Robustness increases with
connectance. Ecology letters 5:558--567.

\bibitem[\citeproctext]{ref-engen2008}
Engen, S., R. Lande, and B.-E. Sæther. 2008. A general model for
analyzing taylor's spatial scaling laws. Ecology 89:2612--2622.

\bibitem[\citeproctext]{ref-etienne2007}
Etienne, R. S. 2007. A neutral sampling formula for multiple samples and
an {``exact''} test of neutrality. Ecology letters 10:608--618.

\bibitem[\citeproctext]{ref-harris2016}
Harris, D. J. 2016. Inferring species interactions from co-occurrence
data with markov networks. Ecology 97:3308--3314.

\bibitem[\citeproctext]{ref-harte2011}
Harte, J. 2011. The maximum entropy theory of ecology. Oxford University
Press.

\bibitem[\citeproctext]{ref-holt1994}
Holt, R. D., and J. Lawton. 1994. The ecological consequences of shared
natural enemies. Annual review of Ecology and Systematics:495--520.

\bibitem[\citeproctext]{ref-hubbell2001}
Hubbell, S. P. 2001. The unified neutral theory of biodiversity and
biogeography. Princeton University Press.

\bibitem[\citeproctext]{ref-hutchinson1961}
Hutchinson, G. E. 1961. The paradox of the plankton. The American
Naturalist 95:137--145.

\bibitem[\citeproctext]{ref-johnson2019}
Johnson, C. A., and J. L. Bronstein. 2019. Coexistence and competitive
exclusion in mutualism. Wiley Online Library.

\bibitem[\citeproctext]{ref-kendall1949}
Kendall, D. G. 1949. Stochastic processes and population growth. Journal
of the Royal Statistical Society. Series B (Methodological) 11:230--282.

\bibitem[\citeproctext]{ref-kurtz2015}
Kurtz, Z. D., C. L. Müller, E. R. Miraldi, D. R. Littman, M. J. Blaser,
and R. A. Bonneau. 2015. Sparse and compositionally robust inference of
microbial ecological networks. PLoS computational biology 11:e1004226.

\bibitem[\citeproctext]{ref-ladau2008}
Ladau, J. 2008. Validation of null model tests using neyman--pearson
hypothesis testing theory. Theoretical Ecology 1:241--248.

\bibitem[\citeproctext]{ref-ladau2017}
Ladau, J. 2017. The architecture and design of ecological null models.
bioRxiv:195131.

\bibitem[\citeproctext]{ref-legendre1993}
Legendre, P. 1993. Spatial autocorrelation: Trouble or new paradigm?
Ecology 74:1659--1673.

\bibitem[\citeproctext]{ref-lennon2000}
Lennon, J. J. 2000. Red-shifts and red herrings in geographical ecology.
Ecography 23:101--113.

\bibitem[\citeproctext]{ref-levine2009}
Levine, J. M., and J. HilleRisLambers. 2009. The importance of niches
for the maintenance of species diversity. Nature 461:254--257.

\bibitem[\citeproctext]{ref-losapio2021}
Losapio, G., C. Schöb, P. P. Staniczenko, F. Carrara, G. M. Palamara, C.
M. De Moraes, M. C. Mescher, R. W. Brooker, B. J. Butterfield, R. M.
Callaway, and others. 2021. Network motifs involving both competition
and facilitation predict biodiversity in alpine plant communities.
Proceedings of the National Academy of Sciences 118:e2005759118.

\bibitem[\citeproctext]{ref-may1972}
May, R. M. 1972. Will a large complex system be stable? Nature
238:413--414.

\bibitem[\citeproctext]{ref-mcgill2003}
McGill, B., and C. Collins. 2003. A unified theory for macroecology
based on spatial patterns of abundance. Evolutionary Ecology Research
5:469--492.

\bibitem[\citeproctext]{ref-morueta2016}
Morueta-Holme, N., B. Blonder, B. Sandel, B. J. McGill, R. K. Peet, J.
E. Ott, C. Violle, B. J. Enquist, P. M. Jørgensen, and J.-C. Svenning.
2016. A network approach for inferring species associations from
co-occurrence data. Ecography 39:1139--1150.

\bibitem[\citeproctext]{ref-proulx2005}
Proulx, S. R., D. E. Promislow, and P. C. Phillips. 2005. Network
thinking in ecology and evolution. Trends in ecology \& evolution
20:345--353.

\bibitem[\citeproctext]{ref-taylor1961}
Taylor, L. R. 1961. Aggregation, variance and the mean. Nature
189:732--735.

\bibitem[\citeproctext]{ref-ulrich2010}
Ulrich, W., and N. J. Gotelli. 2010. Null model analysis of species
associations using abundance data. Ecology 91:3384--3397.

\bibitem[\citeproctext]{ref-williams2000}
Williams, R. J., and N. D. Martinez. 2000. Simple rules yield complex
food webs. Nature 404:180--183.

\bibitem[\citeproctext]{ref-zhou2010}
Zhou, J., Y. Deng, F. Luo, Z. He, Q. Tu, X. Zhi, J. Zhou, Y. Deng, F.
Luo, Z. He, and others. 2010. Functional molecular ecological networks.
mBio 1:e00169--10.

\bibitem[\citeproctext]{ref-zillio2010}
Zillio, T., and F. He. 2010. Modeling spatial aggregation of finite
populations. Ecology 91:3698--3706.

\end{CSLReferences}




\end{document}
